\chapter[Zoiper instalazioa]{Zoiper instalazioa}
\label{chp:zoiper-install}

\localtableofcontents

\section{Sarrera}

Zoiper softwarearekin gure asterisk zerbitzariarentzat bezeroak sortu ditzakegu. Hainbat bezero ahal izateko ordenagailu bakar batean, instalazioa Docker irudiak erabiliz egingo dugu.

\section{Instalazioa}

Hurrengo pausoak jarraituz Zoiper aplikazioa instalatuko da Docker edukiontzi batean.
\begin{enumerate}
    \item Karpeta prestatu instalazioa egiteko: \mintinline{text}{~/Documents/VoIP/Zoiper/}
    \begin{minted}{console}
    $ mkdir -p ~/Documents/VoIP/Zoiper
    \end{minted}
    \item Zoiper-en kodea deskargatu aurretik sortutako karpetan.
    \begin{important}
        Baliteke web orriak cookiak eskatzea deskarga egiteko. Horrela bada, nabigatzaile batetik deskargatu beharko da.
    \end{important}
    \begin{minted}{console}
    $ cd ~/Documents/VoIP/Zoiper
    $ wget -O Zoiper https://www.zoiper.com/en/voip-softphone/download/zoiper5/for/linux
    \end{minted}
    \item Fitxategia deskonprimitu.
    \begin{minted}{console}
    $ tar -Jxvf Zoiper
    \end{minted}
    \item Docker edukiontzia sortzeko konfigurazio fitxategia sortu: \mintinline{text}{~/Documents/VoIP/Zoiper/Dockerfile}
    \begin{minted}{docker}
    FROM ubuntu:latest
    
    RUN apt-get update && apt-get install -y libnotify4 libxss1 libv4l-dev libpulse-dev alsa-utils libgtk2.0-0 libnss3 libatk-bridge2.0-0 libdrm-dev libgbm-dev libxkbcommon-dev libasound-dev libasound2-plugins iproute2 libgl1 nano vim
    \end{minted}
    \item Docker irudia sortu hurrengoa exekutatuz.
    \begin{minted}{console}
    $ docker build -t ubuntu:zoiper .
    \end{minted}
\end{enumerate}

\section{Exekuzioa}

Hurrengo pausoek Zoiper exekutatzeko pausoak adierazten ditu.
\begin{enumerate}
    \item Lehenengo, X11 aplikazioak exekutatzeko ahalmena eman Linuxeko erabiltzaile guztiei:
    \begin{minted}{console}
    $ xhost +Local:*
    \end{minted}
    \item Docker edukiontzia exekutatu hurrengo agindua erabiliz:
    \begin{minted}{console}
    $ docker run -it -e DISPLAY \
    -v /tmp/.X11-unix/:/tmp/.X11-unix \
    -v /home/$USER/Documents/VoIP/Zoiper:/Zoiper \
    -v /etc/alsa:/etc/alsa \
    -v /usr/share/alsa:/usr/share/alsa \
    -v ~/.config/pulse:/.config/pulse \
    -v "/run/user/$UID/pulse/native:/run/user/$UID/pulse/native" \
    --env "PULSE_SERVER=unix:/run/user/$UID/pulse/native" \
    --device /dev/dri:/dev/dri \
    --device /dev/snd \
    ubuntu:zoiper
    \end{minted}
    \item Edukiontzi barruan erabiltzaile berri bat sortu, eta erabiltzaile horrekin sesioa hasi.
    \begin{minted}{console}
    # nereErabiltzailea="Erabiltzaile-izena"
    # useradd -u $(ls /run/user) -s /bin/bash -m $nereErabiltzailea
    # su $nereErabiltzailea
    \end{minted}
    \item Zoiper exekutatu.
    \begin{minted}{console}
    $ cd /Zoiper/Zoiper5 && ./zoiper
    \end{minted}
\end{enumerate}